\documentclass[11pt,a4paper]{article}
\usepackage[T1]{fontenc}
\usepackage[utf8]{inputenc}
\usepackage{lmodern}
\usepackage{amsmath,amssymb,amsthm,mathtools}
\usepackage{esint,graphicx}
\usepackage[a4paper,margin=22mm,headheight=14pt]{geometry}
\usepackage{microtype}
\usepackage{booktabs,longtable,array,enumitem}
\newcolumntype{L}[1]{>{\raggedright\arraybackslash}p{#1}}
\usepackage[dvipsnames]{xcolor}
\usepackage{fancyhdr}
\usepackage{hyperref}

\newtheorem{theorem}{Theorem}[section]
\newtheorem{lemma}[theorem]{Lemma}
\newtheorem{proposition}[theorem]{Proposition}
\newtheorem{corollary}[theorem]{Corollary}
\theoremstyle{remark}
\newtheorem{remark}[theorem]{Remark}
\newtheorem{definition}[theorem]{Definition}
\newtheorem{claim}[theorem]{Claim}
\setlist{nosep,leftmargin=1.5em}
\setlength{\emergencystretch}{2em}
\pagestyle{fancy}
\fancyhf{}
\fancyhead[L]{\small AIM 728}
\fancyhead[R]{\small Revised 7 November 2026}
\fancyfoot[C]{\thepage}
\renewcommand{\headrulewidth}{0.3pt}
\hypersetup{colorlinks=true,linkcolor=MidnightBlue,urlcolor=MidnightBlue,citecolor=MidnightBlue,pdftitle={AIM 728},pdfauthor={Matthew J. Colbrook; results found by AI}}
\begin{document}
\title{The weak endpoint for derivatives of polynomial roots}
\author{Matthew J. Colbrook\thanks{The solutions and partial results in this document were found by AI. The supplied research and the separate repository review are AI work; human peer review and formal verification are not claimed.}\\[4pt]
\small Department of Applied Mathematics and Theoretical Physics\\
\small University of Cambridge, Cambridge CB3 0WA, United Kingdom\\
\small\href{mailto:m.colbrook@damtp.cam.ac.uk}{m.colbrook@damtp.cam.ac.uk}}
\date{7 November 2026}
\maketitle
\begin{abstract}
Continuous roots of a $C^{n-1,1}$ polynomial curve have derivative in weak $L^{n/(n-1)}$, with a coefficient bound. The conclusion concerns exactly the catalogue target; no historical priority claim is made.
\end{abstract}
\noindent\textbf{Catalogue scope.} AIM 728; complete resolution. \href{https://github.com/MColbrook/AIM/blob/1e564633494c109edfd3ec46da1366adcff925a5/problems/728-do-continuous-polynomial-roots-have-weak-endpoint-derivative-regularity.md}{Canonical target at revision \texttt{1e56463}}.


\section{Statement and argument}
\label{sec:728}

The answer to the stated question is affirmative.  The argument below combines
known estimates for radicals and lower-degree polynomials with a maximal-function
covering argument.  In particular, it does not sum local weak-Lebesgue
quasinorms, whose failure of additivity is the obstruction mentioned in
\cite[Open Problem 2]{ParusinskiRainer}.

For an interval $I$, write
\[
 \|f\|_{C^{m,1}(I)}
 =\max\left\{\max_{0\le s\le m}\|f^{(s)}\|_{L^\infty(I)},
                 \operatorname{Lip}_I(f^{(m)})\right\},
 \qquad
 \|f\|_{p,\infty;I}
 =\sup_{u>0}u\,|\{|f|>u\}\cap I|^{1/p}.
\]
Equivalent conventions for the first norm only change constants.

\begin{theorem}
\label{thm:728}
For every integer $n\ge3$ there is a constant $C_n$ with the following
property.  Let
\[
 P_t(z)=z^n+\sum_{j=1}^n a_j(t)z^{n-j},\qquad
 a_j\in C^{n-1,1}([0,1];\mathbb C),
\]
and let $\lambda:(0,1)\to\mathbb C$ be any continuous root.  Then, with
$p=n/(n-1)$,
\[
 \|\lambda'\|_{p,\infty;(0,1)}
 \le C_n\max_{1\le j\le n}
               \|a_j\|_{C^{n-1,1}([0,1])}^{1/j}.
 \tag{728.1}
\]
Consequently every root in the exact hypotheses of AIM 728 has derivative in
$L^{n/(n-1),\infty}(I)$.
\end{theorem}

\subsection{Established inputs and an elementary interpolation estimate}

We use two previously established results, with their roles kept distinct.
First, the complex radical estimate
\cite[Proposition 1 and Corollary 2]{ParusinskiRainer}, based on
\cite{GhisiGobbino}, says that if $a\in C^{j-1,1}(T;\mathbb C)$, then every
continuous $j$th root is absolutely continuous and
\[
 g_a(t)=
 \begin{cases}
 \displaystyle\frac{|a'(t)|}{j|a(t)|^{1-1/j}},&a(t)\ne0,\\[4pt]
 0,&a(t)=0
 \end{cases}
 \quad\hbox{belongs to } L^{j/(j-1),\infty}(T).
 \tag{728.2}
\]
Its weak norm is bounded by a constant depending only on $j$ times
\[
 \max\left\{
 \operatorname{Lip}_T(a^{(j-1)})^{1/j}|T|^{(j-1)/j},
 \|a'\|_{L^\infty(T)}^{1/j}\right\}.
\]
The function in (728.2) is the absolute value of the derivative of any
continuous root, almost everywhere.

Second, the polynomial-root theorem
\cite[Theorem 1]{ParusinskiRainer} gives a uniform strong $L^q$ derivative
estimate in degree $d$ whenever $q<d/(d-1)$, in terms of the
$C^{d-1,1}$ coefficient norms and the interval length.  We will apply this
estimate only with $d<n$.  Degree one is immediate.  We also use the already
known absolute continuity of degree-$n$ roots, which does not assume the
endpoint conclusion.

\begin{lemma}
\label{lem:728interpolation}
If $f\in C^{m-1,1}(I)$ and $\ell=|I|>0$, then
\[
 \ell^s\|f^{(s)}\|_{L^\infty(I)}
 \le C_m\left(\|f\|_{L^\infty(I)}
           +\ell^m\operatorname{Lip}_I(f^{(m-1)})\right),
 \qquad 0\le s\le m-1.
 \tag{728.3}
\]
\end{lemma}
\begin{proof}
Scale $I$ to a unit interval and subtract the Taylor polynomial of degree
$m-1$ at its midpoint.  Lipschitz continuity of the top derivative bounds
all derivatives of the remainder through order $m-1$ by
$C_m\operatorname{Lip}(f^{(m-1)})$.  On the finite-dimensional space of
polynomials of degree at most $m-1$, every derivative supremum is bounded
by a constant times the polynomial's supremum.  The latter is at most
$\|f\|_\infty+C_m\operatorname{Lip}(f^{(m-1)})$.  Rescaling proves
(728.3).  The argument applies to complex-valued functions as well.
\end{proof}

\subsection{Reduction and the stopping intervals}

Put
\[
 R=\max_{1\le j\le n}\|a_j\|_{C^{n-1,1}([0,1])}^{1/j}.
\]
The Tschirnhausen substitution $z=w-a_1(t)/n$ removes the coefficient of
$w^{n-1}$.  Its new coefficients $\widetilde a_j$, $2\le j\le n$, are
polynomials of weighted degree $j$ in the original coefficients, where
$a_i$ has weight $i$.  Leibniz's rule therefore gives
\[
 \|\widetilde a_j\|_{C^{n-1,1}([0,1])}\le C_nR^j.
\]
Extend these functions to $T=(-2,3)$ by a bounded $C^{n-1,1}$ extension
operator.  Endpoint Taylor extensions already provide such an operator on
this fixed larger interval.  After changing $C_n$, the same bounds hold
on $T$.  Since $\widetilde\lambda=\lambda+a_1/n$, it suffices to prove the
estimate for the transformed root and then add the bounded derivative
$a_1'/n$.  Suppress tildes for the moment, so that $a_1=0$.

Define
\[
 A_0=\max_{2\le j\le n}\|a_j\|_{C^{n-1,1}(T)}^{1/j},
 \qquad A(t)=\max_{2\le j\le n}|a_j(t)|^{1/j}.
\]
If $A_0=0$, the transformed root is zero.  Hence assume $A_0>0$.  Let $g_j$
be the function in (728.2) for $a_j$, and set
\[
 h(t)=2A_0+\sum_{j=2}^n g_j(t)\quad(t\in T),
 \qquad h(t)=0\quad(t\notin T).
 \tag{728.4}
\]
For $j<n$, the $C^{j-1,1}$ norm needed in (728.2) is controlled by the
$C^{n-1,1}$ norm; for $j=n$ it is exactly the available regularity.
Since $T$ has fixed finite length and $j/(j-1)\ge p$, the radical estimate
and the finite-sum inequality for weak norms give
\[
 h\in L^{p,\infty}(\mathbb R),\qquad
 \|h\|_{p,\infty;\mathbb R}\le C_nA_0.
 \tag{728.5}
\]
In particular $h$ is locally integrable, since $p>1$.

Fix a sufficiently small constant $\eta>0$, depending only on $n$; its
required smallness will be specified by the splitting construction below.
Take $\eta<1/3$.  For each
\[
 t\in E=\{t\in(0,1):A(t)>0\},
\]
choose the centered open interval $I_t$ satisfying
\[
 \int_{I_t}h(s)\,ds=\eta A(t).
 \tag{728.6}
\]
The radius exists uniquely: the integral is continuous and strictly
increasing while the interval lies in $T$, and $h\ge2A_0$ there.  If
$\ell_t=|I_t|$, then
\[
 \ell_t\le\frac{\eta A(t)}{2A_0}\le\frac\eta2.
 \tag{728.7}
\]
Thus $I_t\subset T$, and $I_t\cap(0,1)$ has length between $\ell_t/2$
and $\ell_t$.

Choose continuous $j$th roots of $a_j$ on $I_t$.  Absolute continuity and
(728.6) imply that their variation from the center is at most $\eta A(t)$.
With $A=A(t)$ this gives
\[
 \sup_{I_t}|a_j|\le(1+\eta)^jA^j,
 \qquad
 \sup_{s\in I_t}\frac{|a_j(s)-a_j(t)|}{A^j}
       \le j(1+\eta)^{j-1}\eta.
 \tag{728.8}
\]
Let $L_j=\operatorname{Lip}_T(a_j^{(n-1)})\le A_0^j$.  From (728.7),
\[
 \frac{L_j\ell_t^n}{A^j}
 \le \left(\frac\eta2\right)^n
           \left(\frac{A}{A_0}\right)^{n-j}
 \le \left(\frac\eta2\right)^n.
 \tag{728.9}
\]
Apply Lemma~\ref{lem:728interpolation} with $m=n$.  Equations
(728.8)--(728.9) show that all the functions
\[
 s\longmapsto\frac{a_j(t+\ell_ts)}{A^j},
 \qquad -\tfrac12\le s\le\tfrac12,
 \tag{728.10}
\]
have uniformly bounded $C^{n-1,1}$ norms, with constants depending only
on $n$.

\subsection{Uniform splitting and the gain in the local exponent}

Consider the compact set
\[
 K=\left\{(c_2,\ldots,c_n)\in\mathbb C^{n-1}:
                   \max_{2\le j\le n}|c_j|^{1/j}=1\right\}.
\]
For every $c\in K$, the polynomial
$z^n+\sum_{j=2}^nc_jz^{n-j}$ has at least two distinct roots.  Indeed, if
all roots were equal, the missing $z^{n-1}$ coefficient would force their
common value to be zero, contradicting $c\in K$.

Group the distinct roots into two nonempty groups and form the corresponding
monic coprime factors $B$ and $C$.  The derivative of the multiplication map
on their coefficient vectors is
\[
 (\dot B,\dot C)\longmapsto C\dot B+B\dot C,
 \qquad \deg\dot B<\deg B,\quad\deg\dot C<\deg C.
\]
This map is injective: coprimality implies that $B$ divides $\dot B$ when
$C\dot B+B\dot C=0$, and the degree bound forces $\dot B=0$, then
$\dot C=0$.  Equal dimensions make it invertible.  The analytic implicit
function theorem therefore gives an analytic factorization chart near
$c$, in which both factors remain coprime and have degrees less than $n$.

Choose nested neighborhoods for these charts, with the closures of the
inner neighborhoods contained in the outer ones.  Compactness of $K$
provides a finite inner cover, a positive uniform margin, and uniform bounds
on the finitely many derivatives of the factor-coefficient maps needed
below.  At the center of $I_t$, the fixed-scale coefficient vector
\[
 \left(\frac{a_2(t)}{A^2},\ldots,\frac{a_n(t)}{A^n}\right)
\]
belongs to $K$.  Choose $\eta$ small enough that (728.8) keeps its entire
path on $I_t$ inside a single outer chart, within the chosen margin.
Together with (728.10), the chain rule then gives uniformly bounded
$C^{n-1,1}$ norms for the coefficients of both normalized factors on the
unit rescaled interval.

Fix
\[
 \frac{n}{n-1}<q<\frac{n-1}{n-2}.
 \tag{728.11}
\]
The root under consideration belongs throughout the connected interval
$I_t\cap(0,1)$ to one of these two factors.  It cannot switch factors:
their root sets are disjoint at every parameter, so the choice of factor
is locally constant.  If this factor has degree $d<n$, then
$q<d/(d-1)$ when $d\ge2$; degree one is elementary.  Apply the established
strong subcritical theorem in degree $d$ to
\[
 \widehat\lambda(s)=\frac{\lambda(t+\ell_ts)}{A}.
\]
The rescaled parameter interval has length between $1/2$ and $1$, and
its normalized factor coefficients have bounded $C^{d-1,1}$ norms.
Consequently
\[
 \int_{I_t\cap(0,1)}|\lambda'(s)|^q\,ds
       \le C_{n,q} A(t)^q |I_t|^{1-q}.
 \tag{728.12}
\]
This is the essential gain: the local exponent is strictly greater than
$p$, because the polynomial has split into smaller degrees.  No
above-endpoint estimate in degree $n$ has been assumed.

\subsection{Patching by disjoint cores and a maximal function}

Let $J_t$ be the concentric interval with one fifth the length of $I_t$.
The $5r$ covering lemma selects a countable pairwise disjoint family
$J_i=J_{t_i}$ such that
\[
 E\subset\bigcup_i5J_i=\bigcup_iI_{t_i}.
\]
Write $I_i=I_{t_i}$ and $A_i=A(t_i)$.  We do not need bounded overlap of
the larger intervals $I_i$.

Let $\mathcal M$ be the uncentered Hardy--Littlewood maximal operator on
$\mathbb R$.  By Chebyshev's inequality and (728.12), for $u>0$,
\[
 \big|\{s\in I_i\cap(0,1):|\lambda'(s)|>u\}\big|
 \le |I_i|\min\left\{1,
       C_{n,q}\left(\frac{A_i}{u|I_i|}\right)^q\right\}.
\]
Equation (728.6) identifies
$A_i/|I_i|=\eta^{-1}|I_i|^{-1}\int_{I_i}h$.
For every $x\in J_i\subset I_i$, this average is at most
$\mathcal Mh(x)$.  Since $|I_i|=5|J_i|$, summing over the disjoint cores
therefore gives
\[
 \big|\{s\in E:|\lambda'(s)|>u\}\big|
 \le5\int_{\mathbb R}
           \min\{1,C_{n,q}(\mathcal Mh(x)/u)^q\}\,dx.
 \tag{728.13}
\]
The known absolute continuity of the root implies that its derivative
vanishes almost everywhere on $\{A=0\}$, where the transformed root is
zero.  Thus the left side of (728.13) can be replaced by the distribution
function on all of $(0,1)$.

For completeness, $\mathcal M$ is bounded on $L^{p,\infty}$ when $p>1$.
One can see this directly by splitting $h$ at height $u/2$: the maximal
function of the lower part is at most $u/2$, while the weak $(1,1)$ estimate
applied to the upper part and integration of its distribution function
 give
\[
 |\{\mathcal Mh>u\}|\le C_p u^{-p}\|h\|_{p,\infty}^p.
\]
Put $H=\|\mathcal Mh\|_{p,\infty}$.  For any fixed $D>0$, layer-cake
integration and $q>p$ yield
\[
 \begin{split}
 \int_{\mathbb R}\min\{1,(D\mathcal Mh/u)^q\}\,dx
 &=q(D/u)^q\int_0^{u/D}
                v^{q-1}|\{\mathcal Mh>v\}|\,dv\\
 &\le\frac{q}{q-p}(DH/u)^p.
 \end{split}
 \tag{728.14}
\]
Combining (728.5), (728.13), and (728.14), with a fixed choice of $q$ in
(728.11), proves
\[
 \|\widetilde\lambda'\|_{p,\infty;(0,1)}\le C_nA_0\le C_nR.
\]
Restoring $\lambda=\widetilde\lambda-a_1/n$ and using the finite-sum
inequality for weak norms proves (728.1).

\subsection{An explicit estimate on any bounded interval}

If $I=(a,b)$ has length $\ell$, define
\[
 N_{j,I}=\max\left\{
      \max_{0\le s\le n-1}\ell^s\|a_j^{(s)}\|_{L^\infty(I)},
      \ell^n\operatorname{Lip}_I(a_j^{(n-1)})\right\}.
\]
Rescaling Theorem~\ref{thm:728} gives
\[
 \boxed{\quad
 \|\lambda'\|_{n/(n-1),\infty;I}
       \le C_n\ell^{-1/n}\max_{1\le j\le n}N_{j,I}^{1/j}.
 \quad}
 \tag{728.15}
\]
The dependence on constant coefficient parts is not asserted to be
optimal.  The estimate is finite under exactly the stated hypotheses and
has the correct homogeneity under scaling the root and the parameter.
The example $z^n-t$ still excludes ordinary strong endpoint integrability;
there is no conflict with that sharp obstruction.

\begingroup\small
\begin{thebibliography}{99}
\bibitem{ParusinskiRainer} A. Parusinski and A. Rainer, \emph{Optimal Sobolev regularity of roots of polynomials}, Ann. Sci. Ecole Norm. Sup. (4) \textbf{51} (2018), 1343--1387. \url{https://math.univ-cotedazur.fr/~parus/publis/ens.pdf}. DOI: \url{https://doi.org/10.24033/asens.2376}.
\bibitem{GhisiGobbino} M. Ghisi and M. Gobbino, \emph{Higher order Glaeser inequalities and optimal regularity of roots of real functions}, Ann. Scuola Norm. Sup. Pisa Cl. Sci. (5) \textbf{12} (2013), 1001--1021. \url{https://www.numdam.org/item/ASNSP_2013_5_12_4_1001_0.pdf}.
\end{thebibliography}
\endgroup
\end{document}
